\documentclass[10pt]{article}
% Math Packages
\usepackage{amsmath, mathtools}
\usepackage{amssymb}
\usepackage{amsthm}
\usepackage{amsfonts}
\usepackage{bbm}
\usepackage{breqn}
\usepackage[margin=1in]{geometry}
\usepackage{graphicx}
\usepackage{tikz}
\usetikzlibrary{arrows.meta}
\usetikzlibrary{calc}
\usepackage{forest}
\usepackage{tikz-qtree}
\graphicspath{ {./images/} }
\usepackage{hyperref}
\usepackage[capitalize]{cleveref}
\usepackage[shortlabels]{enumitem}
\usetikzlibrary{arrows,matrix,positioning}
\usepackage{multicol}

\usepackage{listings}
\usepackage{xcolor}

\crefname{problem}{Problem}{Problems}
\Crefname{problem}{Problem}{Problems}

\lstset{
  language=R,
  basicstyle=\ttfamily\footnotesize,
  keywordstyle=\color{blue},
  commentstyle=\color{gray},
  stringstyle=\color{red!70!black},
  framesep=2pt,
  framexleftmargin=0pt, 
  frame=single,
  breaklines=false,
  showstringspaces=false
}
% for the pipe symbol
\usepackage[T1]{fontenc}

% Citing theorems by name. (source: https://tex.stackexchange.com/questions/109843/cleveref-and-named-theorems)
\makeatletter
\newcommand{\ncref}[1]{\cref{#1}\mynameref{#1}{\csname r@#1\endcsname}}

\def\mynameref#1#2{%
  \begingroup
  \edef\@mytxt{#2}%
  \edef\@mytst{\expandafter\@thirdoffive\@mytxt}%
  \ifx\@mytst\empty\else
  \space(\nameref{#1})\fi
  \endgroup
}
\makeatother

% Colorful Notes
\usepackage{color} \definecolor{Red}{rgb}{1,0,0} \definecolor{Blue}{rgb}{0,0,1}
\definecolor{Purple}{rgb}{.5,0,.5} \def\red{\color{Red}} \def\blue{\color{Blue}}
\def\gray{\color{gray}} \def\purple{\color{Purple}}
\newcommand{\rnote}[1]{{\red [#1]}} % \rnote{foo} gives '[foo]' in red
\newcommand{\pnote}[1]{{\purple [#1]}} % \pnote{foo} gives '[foo]' in purple
\newcommand{\bnote}[1]{{\blue #1}} % \bnote{foo} gives 'foo' in blue
\newcommand{\gnote}[1]{{\gray #1}} % \gnote{foo} gives 'foo' in gray
\newcommand{\Max}[1]{{\purple [#1]}} % \bnote{foo} then 'foo' is blue


% Claim numbering (the counter restarts after each proof environment)
\newcounter{claimcount}
\setcounter{claimcount}{0}
\newenvironment{claim}{\refstepcounter{claimcount}\par\addvspace{\medskipamount}\noindent\textbf{Claim \arabic{claimcount}:}}{}
\usepackage{etoolbox}
\AtBeginEnvironment{proof}{\setcounter{claimcount}{0}}
\newenvironment{claimproof}{\par\addvspace{\medskipamount}\noindent\textit{Proof of Claim  \arabic{claimcount}.}}{\hfill\ensuremath{\qedsymbol} \tiny{Claim}

  \medskip}
% Add claim support to cleverref
\crefname{claimcount}{Claim}{Claims}


% Math Environments
\newtheorem{theorem}{Theorem}
\newtheorem{assumption}[theorem]{Assumption}
\newtheorem{lemma}[theorem]{Lemma}
\newtheorem{proposition}[theorem]{Proposition}
\newtheorem{corollary}[theorem]{Corollary}
\newtheorem{question}[theorem]{Question}
\theoremstyle{definition}
\newtheorem*{definition}{Definition}
\newtheorem{remark}[theorem]{Remark}
\newtheorem{example}[theorem]{Example}
\newtheorem{notation}[theorem]{Notation}
\newtheorem{problem}[theorem]{Problem}

% Matrices and Column Vectors. 
\usepackage{stackengine}
\setstackgap{L}{1.0\normalbaselineskip}
\usepackage{tabstackengine}
\setstackEOL{;}% row separator
\setstackTAB{,}% column separator
\setstacktabbedgap{1ex}% inter-column gap
\setstackgap{L}{1.5\normalbaselineskip}% inter-row baselineskip
\let\nmatrix\bracketMatrixstack  %Usage: \nmatrix{1,2,3\4,5,6}
\newcommand\cv[1]{\setstackEOL{,}\bracketMatrixstack{#1}} %usage: \cv{1,2,3}

% Custom Math Coqmmands
\newcommand{\vt}{\vskip 5mm} % vertical space
\newcommand{\fl}{\noindent\textbf} % first line
\newcommand{\Fl}{\vt\noindent\textbf} % first line with space above
\newcommand{\norm}[1]{\left\lVert#1\right\rVert} % norm
\newcommand{\pnorm}[1]{\left\lVert#1\right\rVert_p} % p-norm
\newcommand{\qnorm}[1]{\left\lVert#1\right\rVert_q} % q-norm
\newcommand{\1}[1]{\textbf{1}_{\left[#1\right]}} % indicator function 
\def\limn{\lim_{n\to\infty}} % shortcut for lim as n-> infinity
\def\sumn{\sum_{n=1}^{\infty}} % shortcut for sum from n=1 to infinity
\def\sumkn{\sum_{k=1}^{n}} % shortcut for sum from k=1 to n
\def\sumin{\sum_{i=1}^{n}} % shortcut for sum from i=1 to n
\def\SAs{\sigma\text{-algebras}} % shortcut for $\sigma$-algebras
\def\SA{\sigma\text{-algebra}} % shortcut for $\sigma$-algebra
\def\Ft{\mathcal{F}_t} % time-indexed sigma-algebra (t)
\def\Fs{\mathcal{F}_s} % time-indexed sigma-algebra (s)
\def\F{\mathcal{F}} % sigma-algebra
\def\G{\mathcal{G}} % sigma-algebra
\def\R{\mathbb{R}} % Real numbers
\def\N{\mathbb{N}} % Natural numbers
\def\Z{\mathbb{Z}} % Integers
\def\E{\mathbb{E}} % Expectation
\def\P{\mathbb{P}} % Probability
\def\Q{\mathbb{Q}} % Q probability
\def\dist{\text{dist}} %Text 'dist' for things like 'dist(x,y)'
\newcommand{\indep}{\perp \!\!\! \perp}  %independence symbol
\def\Var{\mathrm{Var}} % Variance
\def\tr{\mathrm{tr}} % trace

% Brackets and Parentheses
\def\[{\left [}
    \def\]{\right ]}
% \def\({\left (}
%   \def\){\right )}



\usepackage{color}
\definecolor{Red}{rgb}{1,0,0}
\definecolor{Blue}{rgb}{0,0,1}
\definecolor{Purple}{rgb}{.5,0,.5}
\def\red{\color{Red}}
\def\blue{\color{Blue}}
\def\gray{\color{gray}}
\def\purple{\color{Purple}}
\definecolor{RoyalBlue}{cmyk}{1, 0.50, 0, 0}
\newcommand{\dempfcolor}[1]{{\color{RoyalBlue}#1}} 
\newcommand{\demph}[1]{\textcolor{RoyalBlue}{\textbf{\slshape #1}}} % Slanted

\usepackage{emoji}
% comment exactly one of the following line to show / hide the solutions
% \newcommand{\solution}[1]{{\purple #1}} % uncomment to show the solutions
\newcommand{\solution}[1]{} % uncomment to hide the solutions




\begin{document}
\begin{center}
  \section*{Math 472: Homework 06}
  \textit{Due Wednesday, March 11}
\end{center}


\begin{problem}[Exercise 8.50]
  Refer to example 8.8 in the textbook. In this example, $p_{1}$ and $p_{2}$
  were used to denote the proportions of refridgerators of brands $A$ and $B$,
  respectively, that failed during the guarantee periods.
  \begin{enumerate}[(a)]
    \item At the approximate 98\% confidence level, what is the largest
    ``believable value'' for the difference in the proportions of failures for
    refridgerators of brands $A$ and $B$?
    \item At the approximate 98\% confidence level, what is the smallest
    ``believable value'' for the difference in the proportions of failures for
    refrigerators of brands $A$ and $B$? 
    \item If $p_1-p_2$ actually equals $0.2251$, which brand has the larger
    proportion of failures during the warranty period? How much larger?
    \item If $p_1-p_2$ actually equals $-0.1451$, which brand has the larger
    proportion of failures during the warranty period larger? How much larger?
    \item As observed in Example 8.8, zero is a believable value of the
    difference. Would you conclude that there is evidence of a difference in the
    proportions of failures (within the warranty period) for the two brands of
    refrigerators? Why?
  \end{enumerate}
\end{problem}


\begin{problem}[Exercise 8.70]
  Let $Y$ be a binomial random variable with parameter $p$. Find the sample
  size necessary to estimate $p$ within $.05$ with probability $.95$ in the
  following situations:
  \begin{itemize}
    \item If $p$ is thought to be approximately $.9$
    \item If no information is known about $p$ (use $p=.5$ in estimating the
    variance of $\widehat{p}$).
  \end{itemize}
\end{problem}

\begin{problem}[Exercise 8.81]
  The carapace lengths (in mm) of ten lobsters examined in a study of the the
  infestation of the \textit{Thenus orientalis} lobster by two types of
  barnacles, \textit{Octolasmis tridens} and \textit{Octolasmis lowei}, are as
  follows:
  \begin{center}
    78 \quad 66 \quad 65 \quad 63 \quad 60 \quad 60 \quad 58 \quad 56 \quad 52
    \quad 50.
  \end{center}
  Find a $95\%$ confidence interval for the mean carapace length (in mm) of
  \textit{Thenus orientalis} lobsters.
\end{problem}

\begin{problem}[Exercise 8.85]
  Two new drugs were given to patients with hypertension. The first drug
  lowered the blood pressure of 16 patients an average of 11 points, with a
  standard deviation of 6 points. The second drug lowered the blood pressure
  of 20 other patients an average of 12 points, with a standard deviation of 8
  points. Determine a 95\% confidence interval for the difference in the
  mean reductions in blood pressure, assuming that the measurements are
  normally distributed with equal variances.
\end{problem}

\begin{problem}[Exercise 9.1]
  Suppose that $Y_{1},Y_{2},Y_{3}$ are a random sample from an exponential
  distribution with density function
  \begin{equation*}
    f(y) = \frac{1}{\theta}e^{-y/\theta}\mathbf{1}_{\left[ y>0 \right] }
  \end{equation*}
  Consider the following estimators
  \begin{equation*}
    \hat{\theta}_{1} =Y_{1}, \quad\quad \widehat{\theta}_{2}= \frac{Y_{1}+Y_{2}}{2}, \quad\quad \widehat{\theta}_{3} = \frac{Y_{1}+2Y_{2}}{3}, \quad\quad \widehat{\theta}_{4} = \min(Y_{1},Y_{2},Y_{3}), \quad\quad \widehat{\theta}_{5}= \overline{Y}
  \end{equation*}
  In problem \#8 of homework 4, you showed that
  $\widehat{\theta}_{1},\widehat{\theta}_{2},\widehat{\theta}_{3}$ and $\widehat{\theta}_{5}$ are
  unbiased estimators of $\theta$. 
  \begin{enumerate}[(a)]
    \item Find $\text{eff}(\widehat{\theta}_{1},\widehat{\theta}_{5})$.
    \item Find $\text{eff}(\widehat{\theta}_{2},\widehat{\theta}_{5})$.
    \item Find $\text{eff}(\widehat{\theta}_{3},\widehat{\theta}_{5})$.
  \end{enumerate}
\end{problem}

\begin{problem}[Exercise 9.6]
  Suppose $Y_{1},\ldots,Y_{n}$ is a random sample of size $n$ from a Poisson
  distribution with mean $\lambda$. Consider
  $\widehat{\lambda}_{1}= \frac{Y_{1}+Y_{2}}{2}$ and
  $\widehat{\lambda}_{2}=\overline{Y}$. Derive the efficiency of
  $\widehat{\lambda}_{1}$ relative to $\widehat{\lambda}_{2}$.
\end{problem}

\begin{problem}[Exercise 9.8]
  Let  $Y_{1},\ldots,Y_{n}$ denote a random sample from a probability density
  function $f(y)$, which has unknown parameter $\theta$. If $\widehat{\theta}$
  is an unbiased estimator of $\theta$, then under very general conditions
  \begin{equation*}
    \Var(\widehat{\theta}) \geq \frac{1}{\mathcal{I}(\theta)}
  \end{equation*}
  where
  \begin{equation*}
    \mathcal{I}(\theta)= n \E\left[- \frac{\partial^{2}}{\partial \theta^{2} }\Big[\log(f(Y)) \Big]  \right].
  \end{equation*}
  (This is known as the Cramer-Rao inequality.) If
  $\Var(\widehat{\theta})=\frac{1}{\mathcal{I}(\theta)}$, the estimator
  $\widehat{\theta}$ is said to be \textit{efficient}.

  Suppose that $f(y)$ is the normal density with mean $\mu$ and variance
  $\sigma^{2}$. Show that $\overline{Y}$ is an efficient estimator of $\mu$.
\end{problem}

\begin{problem}[Exercise 9.17]
  Suppose that $X_{1},\ldots,X_{n}$ and $Y_{1},\ldots,Y_{n}$ are independent
  random samples from populations with means $\mu_{1}$ and $\mu_{2}$ and
  variances $\sigma^{2}_{1}$ and $\sigma_{2}^{2}$ respectively. Show that
  $\overline{X}-\overline{Y}$ is a consistent estimator of $\mu_{1}-\mu_{2}$.
\end{problem}

\begin{problem}[Exercise 9.24]
  Let $Y_{1},\ldots,Y_{n}$ be independent standard normal random variables.
  \begin{enumerate}[(a)]
    \item What is the distribution of $\sum_{i=1}^{n} Y_{i}^{2n}$?
    \item Let $W_{n}= \frac{1}{n} \sum_{i=1}^{n} Y_{i}^{2}$. Does $W_{n}$
    converge in probability to some constant? If so, what is the value of the
    constant?
  \end{enumerate}
\end{problem}


\begin{problem}[Exercise 9.39]
  Let $X_{1},\ldots,X_{n}$ denote a random sample from a Poisson distribution
  with parameter $\lambda$.  Show that $\sum_{i=1}^{n} Y_{i}$ is sufficient
  for $\lambda$. 
\end{problem}




\end{document}